\scanpage{11}
\mspage{4}
% letterhead
\begin{center}
ACADEMIE D'AIX-MARSEILLE\\
\rule{4em}{0.4pt}\\
INSPECTION ACADEMIQUE\\
DE VAUCLUSE\\
\rule{1.6em}{0.4pt}
\end{center}
% ll. 1–2
Mais je conjecture que l'on doit avoir ceci :
% ll. 3–9
\uline{Lemme} (?) Soient $\pi\in\operatorname{Ob}\widetilde\Delta^\wedge$, $Z=(\widetilde\Delta_{/\pi})^{\mathrm o}$ (donc $Z^{\mathrm o}$ une catégorie étale sur $\widetilde\Delta$ \deleted{\unsure{)|$^{\mathrm o}$}}, i.e. $Z$ coétale sur $\widetilde\Delta^{\mathrm o}$), et supposons que \inserted{a)} tout objet 0-connexe de $Z^\wedge$ soit 1-connexe et \inserted{b)} $\pi(\widetilde\Delta_{\deleted{\unsure{n|0}}0})$ est ponctuel, alors $\pi\simeq e_{\widetilde\Delta^\wedge}$ i.e. \inserted{$\pi_n=$}$\pi(\widetilde\Delta_n)$ est uniponctuel pour tout $n\in\unsure{\mathbb N|\mathcal N}$.
% ll. 10–16
Vu ce qui précède, cela implique

\uline{Corollaire 1} Soit $X$ un objet de $\mathrm{Cat}$ tel que tout objet de $X^\wedge$ qui est 0-connexe soit 1-connexe. Alors \deleted{$\unsure{X|Y}$} pour tt objet 0-connexe $F$ de $X^\wedge$, les produits cartésiens $F^n$ $(n\geq1)$ sont 0-connexes, en particulier $F\times F$ est 0-connexe.
% ll. 17–20
Vu \deleted{le fait} que \inserted{l'on sait déjà que} l'hyp. implique que $X^\wedge$ est loc. irréductible (i.e. $X$ satisfait $\unsure{\mathrm{PS}^*1|\mathrm{P5}^*1}$), on \unsure{tr.|trouve|voit} que $X^\wedge$ est loc. tot. 0-connexe, i.e. $X$ \inserted{loc.} \unsure{cofiltrante|cofiltr.te}, donc
% ll. 21–23
\uline{Corollaire 2} L'hyp. du cor 1 implique que $X$ est \inserted{loc.} \unsure{cofiltrante|cofiltr.te|cofiltrant} (et inversement, bien sûr), i.e. les comp. conn. sont \unsure{cofiltrantes|cofiltr.tes}.
% archive sheet number: 10
